% from main.tex
%%%%%%%%%%  scan idx 91  |  folio 78  %%%%%%%%%%
\section*{Exposé I. Généralités sur les conditions de finitude dans les catégories dérivées}
\begin{center}
par L. Illusie
\end{center}

\subsection*{0. Introduction}
L'objet des Exposés I à IV est de développer, avec la généralité qui convient dans ce Séminaire, le formalisme des conditions de finitude dans les catégories dérivées des topos annelés.

Comme il a été dit dans Exp.\ 0, c'est le besoin de définir sur des schémas arbitraires des « groupes de Grothendieck » possédant de bonnes propriétés de variance qui a obligé à généraliser et à assouplir les notions de finitude utilisées jusqu'à présent.

Une notion classique comme celle de faisceau cohérent sur un espace annelé \((X,\mathcal O_X)\) devient sans intérêt dès que \(\mathcal O_X\) n'est plus cohérent. On pourrait songer à la remplacer par la notion de présentation finie, mais celle-ci présente l'inconvénient que le noyau d'un épimorphisme de Modules de présentation finie n'est plus en général de présentation finie. On arrive à une notion satisfaisante en remarquant que, si \(\mathcal O_X\) est cohérent, un faisceau cohérent \(F\) est non seulement de présentation finie mais de \(n\)-présentation finie pour tout \(n\in\mathbb N\), « \(n\)-présentation finie » voulant dire qu'il existe localement une suite exacte
